
\section{\BookRecordIndex{T0222}黎曼映射定理}\label{sec:C-5}

本节利用圆盘上实\BookRecordIndex{T0210}柯西--黎曼算子的\BookRecordIndex{T0221}黎曼--罗赫定理，证明具有光滑边界的单连通域的
两种不同形式的\BookRecordIndex{T0222}黎曼映射定理。全文以
\[
  \mathbb D:=\{z\in\mathbb C\mid |z|\le1\}
\]
表示复平面中的闭单位圆盘，以 $\mathcal J(\mathbb D)$ 表示 $\mathbb D$ 上诱导标准定向的
\BookRecordIndex{T0088}殆复结构空间，以 $\operatorname{Diff}(\mathbb D)$ 表示 $\mathbb D$ 的保定向微分同胚群。

经典形式的\BookRecordIndex{T0222}黎曼映射定理断言：若 $\Omega\subset\mathbb C$ 是非空、不等于 $\mathbb C$ 的
连通单连通开集，则存在全纯微分同胚 $\phi:\operatorname{int}(\mathbb D)\to\Omega$。
如果假设存在从 $\operatorname{int}(\mathbb D)$ 到 $\Omega$ 的某个微分同胚，便可把 $\Omega$
上的标准复结构拉回到 $\operatorname{int}(\mathbb D)$，并用后者上可积的\BookRecordIndex{T0088}殆复结构重新表述
这一结论。在我们的定理版本中，不要求可积性，而假设\BookRecordIndex{T0088}殆复结构延拓到边界。
这一版本的推论包括黎曼映射在边界上的\BookRecordIndex{T0408}正则性，以及二维情形的可积性。
另一方面，我们不处理任意开域。

\begin{theorem}[\BookRecordIndex{T0222}黎曼映射定理]\label{thm:C-5-1}
对每个 $j\in\mathcal J(\mathbb D)$，存在唯一的微分同胚
$\psi\in\operatorname{Diff}(\mathbb D)$，使得
\begin{equation}\label{eq:C-5-1}
  \psi^*i=j,\qquad \psi(0)=0,\qquad \psi(1)=1.
\end{equation}
由此得到的映射 $\mathcal J(\mathbb D)\to\operatorname{Diff}(\mathbb D):j\mapsto\psi$
是连续的。
\end{theorem}

\begin{corollary}\label{cor:C-5-2}
设 $\Omega\subset\mathbb C$ 是具有光滑边界的连通单连通有界开集，则经典\BookRecordIndex{T0222}黎曼映射定理
对 $\Omega$ 成立。此外，每个全纯微分同胚
$\phi:\operatorname{int}(\mathbb D)\to\Omega$ 都延拓为微分同胚
$\mathbb D\to\overline\Omega$。
\end{corollary}

\begin{proof}
任取微分同胚 $f:\mathbb D\to\overline\Omega$，并取定理~\ref{thm:C-5-1} 中对应于
$j:=f^*i$ 的 $\psi:\mathbb D\to\mathbb D$。则
$f\circ\psi^{-1}:\mathbb D\to\overline\Omega$ 是全纯微分同胚。
故它在 $\operatorname{int}(\mathbb D)$ 上的限制与 $\phi$ 只差一个 \BookRecordIndex{T0040}Möbius 变换。
\end{proof}

定理~\ref{thm:C-5-1} 的另一重要推论是：即使流形有边界，其二维复结构仍可积。

\begin{corollary}[可积性]\label{cor:C-5-3}
二维流形 $\Sigma$（可以有边界）上的每个\BookRecordIndex{T0088}殆复结构 $j$ 都可积。
\end{corollary}

\begin{proof}
取 $p_0\in\Sigma$，并在它的开邻域 $U\subset\Sigma$ 上选取坐标卡
$\phi:U\to\Omega$，其中 $\Omega$ 是闭上半平面
$\mathbb H:=\{\zeta\in\mathbb C\mid\operatorname{Im}\zeta\ge0\}$ 中的开集。
将 $\phi$ 与 \BookRecordIndex{T0040}Möbius 变换
\begin{equation}\label{eq:C-5-2}
  \rho:\mathbb H\to\mathbb D,\qquad
  \rho(\zeta):=\frac{i-\zeta}{i+\zeta}
\end{equation}
复合，得到微分同胚 $\rho\circ\phi:U\to\rho(\Omega)\subset\mathbb D$。
必要时缩小 $U$，可不失一般性地假设 $\rho_*\phi_*j$ 延拓为
$\mathbb D$ 上的\BookRecordIndex{T0088}殆复结构 $j'$。由定理~\ref{thm:C-5-1}，存在微分同胚
$\psi:\mathbb D\to\mathbb D$ 使 $\psi^*i=j'$。因而
$\phi^*\rho^*\psi^*i=j$，于是
\[
  \rho^{-1}\circ\psi\circ\rho\circ\phi:
  U\longrightarrow\rho^{-1}\circ\psi\circ\rho(\Omega)\subset\mathbb H
\]
是 $U$ 上取值于 $\mathbb H$ 的开子集的全纯坐标卡。这就证明了推论。
\end{proof}

\begin{exercise}\label{ex:C-5-4}
设 $(\Sigma,j)$ 是带边界的黎曼曲面，$\Gamma\subset\partial\Sigma$ 是一个与圆周微分同胚
的边界连通分支。证明：当 $\varepsilon>0$ 充分小时，存在从 $\Gamma$ 的一个邻域
到\BookRecordIndex{T0170}环域 $\{z\in\mathbb C\mid1-\varepsilon\le|z|\le1\}$ 的全纯微分同胚。
\end{exercise}

\begin{corollary}\label{cor:C-5-5}
考虑嵌入
\begin{equation}\label{eq:C-5-3}
  \operatorname{PSL}(2,\mathbb R)\longrightarrow\operatorname{Diff}(\mathbb D),
\end{equation}
它把 $\operatorname{PSL}(2,\mathbb R)$ 的每个元素对应于：在
\eqref{eq:C-5-2} 的微分同胚
$\rho:\mathbb H\cup\{\infty\}\to\mathbb D$ 下与相应的 $\mathbb H$ 的 \BookRecordIndex{T0040}Möbius 变换共轭
所得的映射。映射~\eqref{eq:C-5-3} 是同伦等价。特别地，
$\operatorname{Diff}(\mathbb D)$ 是连通的。
\end{corollary}

\begin{proof}
令 $\mathcal J(\mathbb D)\to\operatorname{Diff}(\mathbb D):j\mapsto\psi_j$
把每个\BookRecordIndex{T0088}殆复结构 $j$ 对应于定理~\ref{thm:C-5-1} 中的微分同胚
$\psi_j:=\psi$。则映射
\[
  \operatorname{PSL}(2,\mathbb R)\times\mathcal J(\mathbb D)
  \longrightarrow\operatorname{Diff}(\mathbb D),\qquad
  (\phi,j)\longmapsto\phi\circ\psi_j
\]
是同胚，其逆映射为
$f\mapsto(f\circ\psi_{f^*i}^{-1},f^*i)$。因此推论~\ref{cor:C-5-5}
由 $\mathcal J(\mathbb D)$ 可缩\BookRecordIndex{T0326}推出。
\end{proof}

\begin{remark}\label{rem:C-5-6}
我们在附录 E 给出定理~\ref{thm:C-5-1} 的局部版本的另一种证明
（见定理 E.3.1）。注意，\BookRecordIndex{T0088}殆复结构 $j\in\mathcal J(\mathbb D)$
对应于唯一的度量共形类，使向量 $\xi$ 与 $j\xi$ 相差四分之一周的旋转
（该共形类由一切与 $j$ \BookRecordIndex{T0348}相容的度量组成）。因此，定理~\ref{thm:C-5-1}
也可看作圆盘的\BookRecordIndex{T0375}一致化定理；参见推论 E.3.2。
\end{remark}

我们关于定理~\ref{thm:C-5-1} 的第一个证明基于 Earle 与 Eells~\srcite{95} 的想法。
它只适用于 $j$ 可积的情形：为建立\BookRecordIndex{T0210}柯西--黎曼算子
$\bar\partial_j:\Omega^0_{T\partial\mathbb D}(\mathbb D,T\mathbb D)
\to\Omega^{0,1}_j(\mathbb D,T\mathbb D)$ 的弗雷德霍姆性质，这一点是必需的。
作为交换，该证明在精神上是几何性的，几乎不使用分析。第二个证明则基于延拓论证，
不要求可积性。

\emph{定理~\ref{thm:C-5-1} 的第一个证明：可积情形。}
依 Earle 与 Eells~\srcite{95}，考虑图式
\[
  \operatorname{PSL}(2,\mathbb R)
  \longrightarrow\operatorname{Diff}(\mathbb D)
  \longrightarrow\mathcal J(\mathbb D),
\]
其中第一个映射是嵌入~\eqref{eq:C-5-3}，第二个映射由
$\psi\mapsto\psi^*i$ 给出。

\emph{第一步：映射 $\operatorname{Diff}(\mathbb D)\to\mathcal J(\mathbb D)$，
$\psi\mapsto\psi^*i$ 是一个\BookRecordIndex{T0279}纤维化。}\sourcefootnote{原书本步先称证明
$\psi\mapsto\psi^*i$ 的路径提升，却在随后的初值与方程中使用
$\psi_t^*j_t=i$，即 $\psi_t$ 为所需拉回映射的逆；
两者经微分同胚群中的取逆对应。此处保留原证明的方程。}
我们须证明该映射具有\BookRecordIndex{T0238}路径提升性质。因此，设
$\mathbb R\to\mathcal J(\mathbb D):t\mapsto j_t$ 是一条光滑\BookRecordIndex{T0088}殆复结构路径。
给定 $\psi_0\in\operatorname{Diff}(\mathbb D)$，满足 $\psi_0^*j_0=i$，
我们须找出一条从 $\psi_0$ 出发的光滑路径
$\mathbb R\to\operatorname{Diff}(\mathbb D):t\mapsto\psi_t$，使得
\begin{equation}\label{eq:C-5-4}
  \psi_t^*j_t=i.
\end{equation}
依 Moser 的方法，寻找光滑向量场族
$\mathbb R\to\operatorname{Vect}(\mathbb D):t\mapsto\xi_t$，使得
\begin{equation}\label{eq:C-5-5}
  \frac{d}{dt}\psi_t=\xi_t\circ\psi_t.
\end{equation}
所得同痕满足~\eqref{eq:C-5-4}，当且仅当
\begin{equation}\label{eq:C-5-6}
  \mathcal L_{\xi_t}j_t+\frac{d}{dt}j_t=0,
  \qquad \xi_t(\partial\mathbb D)\subset T\partial\mathbb D.
\end{equation}
注意，切空间 $T_j\mathcal J(\mathbb D)$ 由所有满足与 $j$ 反交换的光滑映射
$\alpha:\mathbb D\to\mathbb R^{2\times2}$ 构成。可将其视为
$\mathbb D$ 上取值于切丛 $T\mathbb D=\mathbb D\times\mathbb R^2$ 的复反线性一形式空间：
\[
  T_j\mathcal J(\mathbb D)=\Omega^{0,1}_j(\mathbb D,T\mathbb D).
\]

把 $\mathbb D$ 上的向量场 $\xi$ 视为 $T\mathbb D$ 的截面。
则李导数与黎曼曲面 $(\mathbb D,j)$ 的\BookRecordIndex{T0210}柯西--黎曼算子
\sourcefootnote{原书将向量场空间写作 $\operatorname{Vect}(T\mathbb D)$，
但同一显示式将它等同于 $\Omega^0(\mathbb D,T\mathbb D)$；
按截面的定义，此处补为 $\operatorname{Vect}(\mathbb D)$。}
\[
  \bar\partial_j:\operatorname{Vect}(\mathbb D)
  =\Omega^0(\mathbb D,T\mathbb D)
  \longrightarrow\Omega^{0,1}_j(\mathbb D,T\mathbb D)
  =T_j\mathcal J(\mathbb D)
\]
由下式联系：
\begin{align*}
  (\mathcal L_\xi j)\eta
   &= [j\eta,\xi]-j[\eta,\xi]\\
   &= \nabla_\xi(j\eta)-\nabla_{j\eta}\xi
      -j\nabla_\xi\eta+j\nabla_\eta\xi\\
   &= j\nabla_\eta\xi-\nabla_{j\eta}\xi\\
   &=2j\bar\partial_j\xi(\eta).
\end{align*}
这里 $\nabla$ 是相对于某个与 $j$ \BookRecordIndex{T0348}相容的黎曼度量的 $\mathbb D$ 上的
\BookRecordIndex{T0034}Levi--Civita 联络；第三个等号利用 $\nabla j=0$。最后一式中的
$\bar\partial_j$ 是二维流形 $(\mathbb D,j)$ 的（全纯）切丛上的\BookRecordIndex{T0210}柯西--黎曼算子。
因此，按注~\ref{rem:C-1-2} 的记法，
$\bar\partial_j\xi=\frac12(\nabla\xi+j\nabla\xi\circ j)$。
\BookRecordIndex{T0080}丛对 $(T\mathbb D,T\partial\mathbb D)$ 的\BookRecordIndex{T0062}边界马斯洛夫指标为 $2$，
故由定理~\ref{thm:C-1-10}，对每个 $t$，方程~\eqref{eq:C-5-6} 都有解 $\xi_t$。
再加上条件 $\xi_t(0)=\xi_t(1)=0$，解便唯一，且光滑依赖于 $t$。
这就证明了\BookRecordIndex{T0238}路径提升性质。

\emph{第二步：若 $u:\mathbb D\to\mathbb D$ 是全纯微分同胚，且
$u(0)=0$、$u(1)=1$，则 $u=\operatorname{id}$。}
我们使用反射原理。定义 $v:\mathbb D\to\mathbb D$ 为
$v(z):=\overline{u(\bar z)}$。则 $v$ 是全纯微分同胚，满足
$v(0)=0$，且对 $z\in S^1$ 有 $u(1/z)v(z)=1$。因此，可通过
\[
  u(z):=\frac1{v(1/z)},\qquad |z|\ge1
\]
把 $u$ 延拓成二球面的全纯微分同胚。$u$ 的导数形如
\[
  u'(z)=\frac{v'(1/z)}{z^2v(1/z)^2},\qquad |z|\ge1.
\]
当 $z\to\infty$ 时，$zv(1/z)$ 收敛到 $v'(0)\ne0$，故 $u'$
是有界全纯函数，由刘维尔定理它是常数。结合 $u(0)=0$、$u(1)=1$，
得到所称的 $u=\operatorname{id}$。

由第一步得存在性，由第二步得唯一性。连续性的证明留作习题。
这就证明了定理~\ref{thm:C-5-1}，前提是 $\mathbb D$ 上每个\BookRecordIndex{T0088}殆复结构均可积。

定理~\ref{thm:C-5-1} 的第二个证明在技术上比第一个更复杂，却有更重要的推论：
它蕴含二维\BookRecordIndex{T0088}殆复结构的可积性，即使在边界点附近亦然。为完成证明，我们须处理如下
基本问题：给定 $\mathbb C$ 上的\BookRecordIndex{T0088}殆复结构 $j$ 和一个 $(i,j)$-全纯圆盘
$u:(\mathbb D,\partial\mathbb D)\to(\mathbb C,S^1)$，且其边界限制
$u|_{\partial\mathbb D}:\partial\mathbb D\to S^1$ 的度为一，为什么 $u$ 是闭圆盘的
微分同胚？较易证明 $u$ 在开单位圆盘上的限制是微分同胚；更微妙的是证明 $u$ 在边界上
不能有\BookRecordIndex{T0233}临界点。相关问题还有边界点处的\BookRecordIndex{T0333}唯一延拓。下文给出基于
定理~\ref{thm:2-3-5} 中 \BookRecordIndex{T0013}Carleman 相似原理的证明。该技术可推广到更高维。

\begin{lemma}\label{lem:C-5-7}
设 $j$ 是 $\mathbb C$ 上 $C^2$ 类\BookRecordIndex{T0088}殆复结构，$u:\mathbb D\to\mathbb C$
是满足
\begin{equation}\label{eq:C-5-7}
  \partial_su+j(u)\partial_tu=0,\qquad
  u(\partial\mathbb D)\subset S^1,\qquad
  \deg(u|_{\partial\mathbb D}:\partial\mathbb D\to S^1)=1
\end{equation}
的 $C^1$ 函数。则 $u$ 是 $\mathbb D$ 的 $C^1$ 微分同胚。
\end{lemma}

\begin{proof}
对 $\zeta\in\mathbb C\setminus S^1$，以 $\deg(u,\zeta)$ 表示 $u$ 在 $\zeta$ 处的
度。由于 $u|_{\partial\mathbb D}:\partial\mathbb D\to S^1$ 的度为一，
映射
\[
  u_\zeta:\partial\mathbb D\to S^1,\qquad
  u_\zeta(z):=\frac{u(z)-\zeta}{|u(z)-\zeta|}
\]
对每个 $|\zeta|<1$ 的 $\zeta$ 也有度一。因此，当 $|\zeta|<1$ 时，
$\deg(u,\zeta)=\deg(u_\zeta)=1$；当 $|\zeta|>1$ 时，
$\deg(u,\zeta)=\deg(u_\zeta)=0$。\sourcefootnote{原书此处将外部点条件印为 $|\zeta|>0$，与紧前 $|\zeta|<1$ 时度为一的结论矛盾；下文也明确在 $|\zeta|>1$ 时使用零度结论。故改为 $|\zeta|>1$。}
把定理~\ref{thm:2-3-5} 中的 \BookRecordIndex{T0013}Carleman 相似原理用于函数
$z\mapsto u(z_0+z)-\zeta$，可知 $u^{-1}(\zeta)$ 中每个点 $z_0$
对局部度都有正贡献；若 $z_0$ 是\BookRecordIndex{T0233}临界点，贡献还大于一。
故当 $|\zeta|>1$ 时 $u^{-1}(\zeta)=\varnothing$；当 $|\zeta|<1$ 时
$\#u^{-1}(\zeta)=1$，且每个 $|\zeta|<1$ 的点都是 $u$ 的\BookRecordIndex{T0409}正则值。
定理~\ref{thm:2-3-5} 还说明 $u(\operatorname{int}\mathbb D)$ 是开集。
因为 $u(\mathbb D)\subset\mathbb D$，这意味着
$u(\operatorname{int}\mathbb D)=\operatorname{int}\mathbb D$。
因此 $u$ 在开单位圆盘上的限制是 $C^1$ 微分同胚。

余下要证明 $u$ 在边界上没有\BookRecordIndex{T0233}临界点。我们首先归约到 $u$ 有边界不动点、
且 $j=i$ 沿边界成立的情形。使用 $\mathbb H$ 比使用 $\mathbb D$ 更方便。
为此，固定 $z_0=e^{i\theta}\in\partial\mathbb D$，定义
$\rho:\mathbb H\to\mathbb D$ 为
\[
  \rho(\zeta):=z_0\frac{i-\zeta}{i+\zeta},\qquad
  \rho(0)=z_0,\qquad \rho(\mathbb H)=\mathbb D\setminus\{-z_0\};
\]
又定义 \BookRecordIndex{T0040}Möbius 变换
\[
  \psi(z):=i\frac{1-\overline{u(z_0)}z}{1+\overline{u(z_0)}z},
  \qquad\psi(u(z_0))=0,\qquad
  \psi(\mathbb D\setminus\{-u(z_0)\})=\mathbb H,
\]
并记
\[
  j':=\psi_*j,\qquad
  u':=\psi\circ u\circ\rho:U'\to\mathbb C,\qquad
  U':=\{\zeta\in\mathbb H\mid u(\rho(\zeta))\ne-u(z_0)\}.
\]
于是 $j'$ 是 $\mathbb C\setminus\{-i\}$ 上的 $C^2$ \BookRecordIndex{T0088}殆复结构，
$u':U'\to\mathbb C\setminus\{-i\}$ 是 $j'$-全纯的 $C^1$ 函数，
满足 $u'(U'\cap\mathbb R)\subset\mathbb R$ 及 $u'(0)=0$。
选取零点的两个邻域 $\Omega'',\Omega'\subset\mathbb C$ 之间的 $C^2$ 微分同胚
$\chi:\Omega''\to\Omega'$，使得
\begin{equation}\label{eq:C-5-8}
  \chi(\zeta)=\zeta,\qquad
  (\chi^*j')(\zeta)=i\quad(\zeta\in\Omega''\cap\mathbb R).
\end{equation}
其一个显式公式为
\[
  \chi(\zeta):=\operatorname{Re}\zeta
       +(\operatorname{Im}\zeta)j'(\zeta)1.
\]
读者可验证，把 $\chi$ 限制到实轴的一个充分小的邻域后，
它是满足~\eqref{eq:C-5-8} 的 $C^2$ 微分同胚。定义
\[
  u'':=\chi^{-1}\circ u':U'':=(u')^{-1}(\Omega')\to\Omega'',
  \qquad j'':=\chi^*j'.
\]
则 $u''$ 和 $j''$ 满足下述引理~\ref{lem:C-5-8} 的假设：
$j''$ 是 $\Omega''$ 上 $C^1$ 类、在实轴上等于 $i$ 的\BookRecordIndex{T0088}殆复结构，且
$u'':U''\to\Omega''$ 是满足
\[
  \partial_su''+j''(u'')\partial_tu''=0,\qquad
  u''(U''\cap\mathbb R)\subset\mathbb R,\qquad u''(0)=0
\]
的 $C^1$ 函数。
由此存在正整数 $\nu$、实数 $\delta>0$ 与 $c\ne0$，使得当 $\zeta$
趋近零时
\[
  u''(\zeta)=c\zeta^\nu+O(|\zeta|^{\nu+\delta}).
\]
因为 $u$ 把 $\mathbb D$ 映到自身，$u''$ 便把上半平面 $\mathbb H$ 映到自身。
因此 $\nu=1$ 且 $c>0$；特别地，$du''(0)\ne0$，从而
$du(z_0)\ne0$。所以 $u$ 在边界上没有\BookRecordIndex{T0233}临界点，进而是闭单位圆盘
$\mathbb D$ 到自身的 $C^1$ 微分同胚。这就证明了引理~\ref{lem:C-5-7}。
\end{proof}

\begin{lemma}\label{lem:C-5-8}
固定实数 $p>2$，记 $\mathbb H\subset\mathbb C$ 为闭上半平面。
设 $j$ 是零点的某个连通开邻域 $\Omega\subset\mathbb C$ 上的
$W^{1,p}$ 类\BookRecordIndex{T0088}殆复结构，$W\subset\mathbb H$ 是零点的一个连通开邻域，
且 $w:W\to\Omega$ 是满足
\begin{equation}\label{eq:C-5-9}
  \partial_sw+j(w)\partial_tw=0,\qquad
  w(W\cap\mathbb R)\subset\mathbb R,\qquad w(0)=0
\end{equation}
的 $W^{1,p}$ 函数。假设 $j(\zeta)=i$ 对 $\zeta\in\mathbb R$ 成立，
且 $w$ 非常值。则存在正整数 $\nu$ 及非零实数 $c$，使得
\begin{equation}\label{eq:C-5-10}
  w(\zeta)=c\zeta^\nu+O(|\zeta|^{\nu+1-2/p}).
\end{equation}
\end{lemma}

\begin{proof}
把 Schwarz 反射式论证与定理~\ref{thm:2-3-5} 的 \BookRecordIndex{T0013}Carleman 相似原理结合，
即可得到结论；后者的证明本身又用到了单位圆盘上的简易\BookRecordIndex{T0221}黎曼--罗赫定理 C.4.1。
令 $\tau:\mathbb C\to\mathbb C$ 表示复共轭。定义
\[
  U:=W\cup\tau(W)\subset\mathbb C,\qquad
  u:U\to\mathbb C,\qquad
  J:U\to\operatorname{End}_{\mathbb R}(\mathbb C)
\]
为
\[
  u(\zeta):=
  \begin{cases}
    w(\zeta),&\zeta\in W,\\
    \overline{w(\bar\zeta)},&\zeta\in\tau(W),
  \end{cases}
  \qquad
  J(\zeta):=
  \begin{cases}
    j(w(\zeta)),&\zeta\in W,\\
    -\tau\circ j(w(\bar\zeta))\circ\tau,&\zeta\in\tau(W).
  \end{cases}
\]
则 $U\subset\mathbb C$ 是零点的开邻域。由于 $w(\zeta)\in\mathbb R$
且 $j(w(\zeta))\in\operatorname{End}_{\mathbb R}(\mathbb C)$
在 $\zeta\in\mathbb R$ 时是乘以 $i$，可得
$u\in W^{1,p}(U,\mathbb C)$、
$J\in W^{1,p}(U,\operatorname{End}_{\mathbb R}(\mathbb C))$。
此外，$u$ 与 $J$ 满足
\begin{equation}\label{eq:C-5-11}
  \partial_su(\zeta)+J(\zeta)\partial_tu(\zeta)=0,
  \qquad u(0)=0,\quad J(0)=i,
\end{equation}
其中 $\zeta=s+it\in U$。因此由定理~\ref{thm:2-3-5}，
存在常数 $r>0$，使
\[
  B_r:=\{\zeta\in\mathbb C\mid|\zeta|<r\}\subset U,
\]
并存在映射
$\Phi\in W^{1,p}(B_r,\operatorname{End}_{\mathbb R}(\mathbb C))$
及全纯函数 $\sigma:B_r\to\mathbb C$，使得 $\Phi(\zeta)$ 可逆，且
\begin{equation}\label{eq:C-5-12}
  u(\zeta)=\Phi(\zeta)\sigma(\zeta),\qquad
  \sigma(0)=0,\qquad J(\zeta)\Phi(\zeta)=\Phi(\zeta)i
\end{equation}
对所有 $\zeta\in B_r$ 成立。因为 $J(0)$ 是乘以 $i$，
而 $J\Phi=\Phi i$，可知 $\Phi(0)$ 为复线性映射。
故存在非零复数 $a$，使 $\Phi(0):\mathbb C\to\mathbb C$ 是乘以 $a$。
由于 $u=\Phi\sigma$ 非常值且 $\Phi$ 处处可逆，$\sigma$ 也非常值。
由 $\sigma(0)=0$，存在正整数 $\nu$ 与非零复数 $b$，使
\[
  \sigma(\zeta)=b\zeta^\nu+O(|\zeta|^{\nu+1}).
\]
记 $c:=ab\in\mathbb C\setminus\{0\}$。于是由~\eqref{eq:C-5-12}，
对 $\zeta\in B_r\cap\mathbb H$ 有
\begin{align*}
  w(\zeta)=u(\zeta)
   &=ab\zeta^\nu
      +(\Phi(\zeta)-\Phi(0))b\zeta^\nu
      +\Phi(\zeta)(\sigma(\zeta)-b\zeta^\nu)\\
   &=c\zeta^\nu+O(|\zeta|^{\nu+1-2/p}).
\end{align*}
最后一个等号来自\BookRecordIndex{T0318}索伯列夫嵌入定理~\ref{thm:B-1-11}：
$\Phi$ 以指数 $1-2/p$ Hölder 连续，故
$\Phi(\zeta)-\Phi(0)=O(|\zeta|^{1-2/p})$。
因为 $u(U\cap\mathbb R)\subset\mathbb R$，$c$ 必为实数。
这就证明了引理~\ref{lem:C-5-8}。
\end{proof}

\emph{定理~\ref{thm:C-5-1} 的第二个证明：一般情形。}
固定整数 $\ell\ge2$ 和实数 $p>2$。令 $\mathcal J^\ell(\mathbb C)$ 为
$\mathbb C$ 上所有 $C^\ell$ 类、且在半径为二的圆盘外等于 $i$ 的\BookRecordIndex{T0088}殆复结构的空间。
于是 $\mathcal J^\ell(\mathbb C)$ 是 Banach 流形。以
$\operatorname{Diff}^1(\mathbb D)$ 表示 $\mathbb D$ 的 $C^1$ 微分同胚群。定义
\[
  \mathcal M^\ell:=
  \{(j,u)\in\mathcal J^\ell(\mathbb C)
     \times\operatorname{Diff}^1(\mathbb D)
       \mid u^*j=i,\ u(0)=0,\ u(1)=1\}
\]
以及
\[
  \mathcal J^\ell_0:=
  \{j\in\mathcal J^\ell(\mathbb C)
    \mid\text{存在 }u\in\operatorname{Diff}^1(\mathbb D)
           \text{ 使 }(j,u)\in\mathcal M^\ell\}.
\]
我们分六步证明 $\mathcal J^\ell_0=\mathcal J^\ell(\mathbb C)$。

\emph{第一步：$\mathcal J^\ell_0$ 非空。}
因为 $(i,\operatorname{id})\in\mathcal M^\ell$，所以
$i\in\mathcal J^\ell_0$。

\emph{第二步：集合}
\[
  \mathcal U^{1,p}:=
  \left\{u\in W^{1,p}(\mathbb D,\mathbb C)\ \middle|\
    \begin{array}{l}
      u(\partial\mathbb D)\subset S^1,\\
      \deg(u|_{\partial\mathbb D}:\partial\mathbb D\to S^1)=1,\\
      u(0)=0,\ u(1)=1
    \end{array}\right\}
\]
是 $W^{1,p}(\mathbb D,\mathbb C)$ 的光滑 Banach 子流形，
且在 $u\in\mathcal U^{1,p}$ 处的切空间为 Banach 空间
\[
  T_u\mathcal U^{1,p}=
  \{\hat u\in W^{1,p}(\mathbb D,\mathbb C)
     \mid\hat u(z)\in T_{u(z)}S^1\ \forall z\in\partial\mathbb D,
     \ \hat u(0)=\hat u(1)=0\}.
\]
这是标准结果，证明留给读者。

\emph{第三步：若 $(j,u)\in\mathcal J^\ell(\mathbb C)\times\mathcal U^{1,p}$
且 $\partial_su+j(u)\partial_tu=0$，则
$u\in\operatorname{Diff}^1(\mathbb D)$。}
由于 $j$ 是 $C^\ell$ 类\BookRecordIndex{T0088}殆复结构，$\ell\ge2$，而 $u$ 是满足边值问题
\eqref{eq:C-5-7} 的 $W^{1,p}$ 解，$p>2$，由定理~\ref{thm:B-4-1}
知 $u$ 为 $W^{\ell,p}$ 函数。再因 $p>2$，由\BookRecordIndex{T0318}索伯列夫嵌入定理
\ref{thm:B-1-11} 知 $u$ 属 $C^{\ell-1}$ 类。因此第三步由
引理~\ref{lem:C-5-7} \BookRecordIndex{T0326}推出。

\emph{第四步：若 $(j,u)\in\mathcal M^\ell$，则算子}
\[
  D_{j,u}:T_u\mathcal U^{1,p}\to L^p(\mathbb D,\mathbb C),
  \qquad
  D_{j,u}\hat u:=\partial_s\hat u+j(u)\partial_t\hat u
               +(dj(u)\hat u)\partial_tu
\]
\emph{是双射。}
每个 $\hat u\in T_u\mathcal U^{1,p}$ 都形如 $\hat u=du\cdot\xi$，
其中 $\xi\in W^{1,p}(\mathbb D,\mathbb C)$ 满足下述条件。
\sourcefootnote{原书将 $\hat u=du\cdot\xi$ 的零点条件转给 $\xi$ 时，
末项印作 $\xi(1)=1$；由线性切空间条件 $\hat u(1)=0$ 及 $du(1)$ 可逆，
应为 $\xi(1)=0$。下文核的计算亦使用此条件。}
\[
  \xi(z)\in T_zS^1\quad(\forall z\in\partial\mathbb D),
  \qquad \xi(0)=0,\qquad\xi(1)=0.
\]
此外，由 $\partial_su+j(u)\partial_tu=0$ 可知，当 $\hat u=du\cdot\xi$ 时，
\[
  D_{j,u}\hat u=du\cdot(\partial_s\xi+i\partial_t\xi),
  \qquad
  D_{j,u}\hat u=0\quad\Longleftrightarrow\quad
  \partial_s\xi+i\partial_t\xi=0.
\]
因此，由定理 C.4.1 证明中的第二步（取 $k=2$）可知，
$D_{j,u}$ 的核中每个元素都形如
\[
  \hat u(z)=du(z)i(a+bz+\bar a z^2),
  \qquad a\in\mathbb C,\quad b\in\mathbb R.
\]
由于 $\hat u(0)=0$，有 $a=0$；又由 $\hat u(1)=0$，有 $b=0$。
故 $D_{j,u}$ 为单射。定理 C.4.1 还说明 $D_{j,u}$ 的\BookRecordIndex{T0129}弗雷德霍姆指标为零
（因为其定义域中的元素满足 $\hat u(0)=\hat u(1)=0$），
于是它是双射。这就证明了第四步。

\emph{第五步：$\mathcal J^\ell_0$ 是 $\mathcal J^\ell(\mathbb C)$ 的开子集。}
由第二步，$\mathcal J^\ell(\mathbb C)\times\mathcal U^{1,p}$ 是光滑
Banach 流形。定义映射
\[
  \mathcal F:\mathcal J^\ell(\mathbb C)\times\mathcal U^{1,p}
  \longrightarrow L^p(\mathbb D,\mathbb C),\qquad
  \mathcal F(j,u):=\partial_su+j(u)\partial_tu.
\]
此映射可连续微分 $\ell$ 次，在 $(j,u)$ 处的导数是线性算子
$d\mathcal F(j,u):T_j\mathcal J^\ell(\mathbb C)\times
T_u\mathcal U^{1,p}\to L^p(\mathbb D,\mathbb C)$，由
\[
  d\mathcal F(j,u)(\hat j,\hat u)
    =D_{j,u}\hat u+\hat j(u)\partial_tu
\]
给出。由第三步，$\mathcal F$ 的零点集是 $\mathcal M^\ell$。
若 $(j,u)\in\mathcal M^\ell$，则由第四步，算子
$D_{j,u}:T_u\mathcal U^{1,p}\to L^p(\mathbb D,\mathbb C)$ 是双射。
因此，由引理~\ref{lem:A-3-6}，$d\mathcal F(j,u)$ 满射且有\BookRecordIndex{T0392}右逆。
于是由\BookRecordIndex{T0380}隐函数定理~\ref{thm:A-3-3}，$\mathcal M^\ell$ 是
$\mathcal J^\ell(\mathbb C)\times\mathcal U^{1,p}$ 的 $C^\ell$ Banach 子流形。
它在 $(j,u)$ 处的切空间是 Banach 空间
\[
  T_{(j,u)}\mathcal M^\ell
   =\{(\hat j,\hat u)\in T_j\mathcal J^\ell(\mathbb C)
       \times T_u\mathcal U^{1,p}\mid
       D_{j,u}\hat u+\hat j(u)\partial_tu=0\}.
\]
由于 $D_{j,u}$ 是双射，投影
\[
  T_{(j,u)}\mathcal M^\ell\to T_j\mathcal J^\ell(\mathbb C),
  \qquad (\hat j,\hat u)\mapsto\hat j
\]
也是双射。这说明 $C^\ell$ Banach 流形之间的投影
$\mathcal M^\ell\to\mathcal J^\ell(\mathbb C):(j,u)\mapsto j$
是局部微分同胚。因此由逆函数定理~\ref{thm:A-3-1}，
其像 $\mathcal J^\ell_0$ 是 $\mathcal J^\ell(\mathbb C)$ 的开子集。
第五步得证。

\emph{第六步：$\mathcal J^\ell_0$ 是 $\mathcal J^\ell(\mathbb C)$ 的闭子集。}
设 $j^\nu$ 是 $\mathcal J^\ell_0$ 中一列在 $C^\ell$ 拓扑下收敛于
$j\in\mathcal J^\ell(\mathbb C)$ 的元素。对每个 $\nu$，选取
$u^\nu\in\operatorname{Diff}^1(\mathbb D)$，使
$(j^\nu,u^\nu)\in\mathcal M^\ell$。由定理~\ref{thm:B-4-1}，
每个 $u^\nu\in W^{\ell,p}(\mathbb D,\mathbb C)$。
我们声称 $u^\nu$ 的一阶导数一致有界。证得这一点后，
由定理~\ref{thm:B-4-2} 可知 $u^\nu$ 在 $C^\ell$ 拓扑下一致有界，
\sourcefootnote{原书在此由 B.4.2 直接说得到 $C^\ell$ 一致上界；
该定理陈述的是 $C^{\ell-1}$ 逐紧收敛子序列，足以\BookRecordIndex{T0326}推出下文所需的
$C^{\ell-1}$ 收敛。这里保留原文较强的中间断言。}
从而有子序列在 $C^{\ell-1}$ 拓扑下收敛于
$u:\mathbb D\to\mathbb C$。极限 $u$ 属于 $\mathcal U^{1,p}$，且满足
$\partial_su+j(u)\partial_tu=0$。因此由第三步
$(j,u)\in\mathcal M^\ell$，故 $j\in\mathcal J^\ell_0$。

为证明所需的一阶导数上界，我们在 $\mathbb C$ 上选取一列黎曼度量，
由\BookRecordIndex{T0088}殆复结构 $j^\nu$ 与标准体积形式 $\omega:=ds\wedge dt$ 按
$\langle\cdot,\cdot\rangle_{j^\nu}:=\omega(\cdot,j^\nu\cdot)$ 确定。
关于该度量，$u^\nu$ 的能量为
\[
  E(u^\nu)=\int_{\mathbb D}|\partial_su^\nu|^2_{j^\nu}\,ds\,dt
    =\int_{\mathbb D}(u^\nu)^*\omega=\pi.
\]
最后一个等号来自 $u^\nu|_{\partial\mathbb D}:\partial\mathbb D\to S^1$
的度等于一。此外，边界落在 $S^1$ 上的每个 $j$-全纯圆盘的能量
都是 $\pi$ 的整数倍。

现在反证，假设序列 $\sup_{\mathbb D}|du^\nu|$ 无界。
则由定理~\ref{thm:4-6-1}，取子序列后，$u^\nu$ 的一阶导数
在某个非空有限集 $Z\subset\mathbb D$ 的补集的每个紧子集上一致有界，
且对每个 $\zeta\in Z$，存在序列 $z^\nu\in\mathbb D$，满足
$z^\nu\to\zeta$ 且 $c^\nu:=|du^\nu(z^\nu)|\to\infty$。
这里 $\zeta\notin\operatorname{int}(\mathbb D)$；否则，\BookRecordIndex{T0412}重标度序列
$v_\nu(z):=u^\nu(z^\nu+z/c^\nu)$ 再取子序列后，
由定理~\ref{thm:B-4-2} 与定理~\ref{thm:4-1-2} 将收敛至
$\mathbb C$ 中的非恒定 $j$-全纯球面，而这样的球面不存在。
故 $Z\subset\partial\mathbb D$，且由定理~\ref{thm:4-6-1}，
\[
  \lim_{\varepsilon\to0}\lim_{\nu\to\infty}
  E(u^\nu,B_\varepsilon(\zeta)\cap\mathbb D)\ge\pi
  \qquad(\forall\zeta\in Z).
\]
因为 $u^\nu$ 的一阶导数在 $\mathbb D\setminus Z$ 的每个紧子集上一致有界，
定理~\ref{thm:B-4-2} 说明再取子序列后，$u^\nu$ 在
$\mathbb D\setminus Z$ 的每个紧子集上以 $C^{\ell-1}$ 拓扑收敛到
$j$-全纯映射 $u:\mathbb D\setminus Z\to\mathbb C$。
因为总能量 $E(u^\nu)=\pi$ 全部消耗在\BookRecordIndex{T0243}冒泡点处（\BookRecordIndex{T0243}冒泡点因而至多一个），
极限 $u$ 必为常值。然而 $u(0)=0$，且
$u(\partial\mathbb D\setminus Z)\subset S^1$，不可能为常值。
这一矛盾证明 $u^\nu$ 的一阶导数确实一致有界。第六步得证。

由于 $\mathcal J^\ell(\mathbb C)$ 连通，第一、五、六步\BookRecordIndex{T0326}推出
$\mathcal J^\ell_0=\mathcal J^\ell(\mathbb C)$。
现在设 $j$ 是 $\mathbb D$ 上的光滑\BookRecordIndex{T0088}殆复结构。把它延拓为
$\mathbb C$ 上的光滑\BookRecordIndex{T0088}殆复结构，仍记 $j$，并使其在半径为二的圆盘外等于 $i$
（见下文习题~\ref{ex:C-5-9}）。于是 $j\in\mathcal J^\ell(\mathbb C)$，
故由刚证得的结论，$j\in\mathcal J^\ell_0$。因此存在 $C^1$ 微分同胚
$u:\mathbb D\to\mathbb D$，使 $(j,u)\in\mathcal M^\ell$。
由 $j$ 光滑及定理~\ref{thm:B-4-1}，$u$ 光滑。
故 $\psi:=u^{-1}:\mathbb D\to\mathbb D$ 是满足~\eqref{eq:C-5-1} 的微分同胚。
这证明存在性。唯一性由定理~\ref{thm:C-5-1} 第一个证明中的第二步\BookRecordIndex{T0326}推出。
为证明连续性，可利用：对每个 $\ell$，投影
$\mathcal M^\ell\to\mathcal J^\ell(\mathbb C)$ 是 $C^\ell$ Banach 流形
之间的 $C^\ell$ 微分同胚。这完成了定理~\ref{thm:C-5-1} 一般情形的证明。

\begin{exercise}\label{ex:C-5-9}
设 $f:(-\infty,0]\to\mathbb R$ 是光滑函数。本习题说明如何将 $f$
延拓成整个 $\mathbb R$ 上的光滑函数，并使其在 $t\ge1$ 时为零。
选取光滑\BookRecordIndex{T0188}截断函数 $\beta:[0,\infty)\to[0,1]$，满足
\[
  \beta(t)=
  \begin{cases}
    1,&0\le t\le\frac12,\\
    0,&t\ge1.
  \end{cases}
\]
对 $k=0,1,2,3,\ldots$，选取 $a_k\in\mathbb R$ 与 $\varepsilon_k>0$，使
\sourcefootnote{原书以下上界式的分母写为 $|a_k|k^k$；若 $a_k=0$，
或在 $k=0$ 时不预设 $0^0$ 的约定，该表达式形式上未定义。
此处照录原式，不替作者增补约定。}
\[
  a_k:=\frac{f^{(k)}(0)}{k!},\qquad
  \varepsilon_k\le\frac1{|a_k|k^k},\qquad
  \varepsilon_k\le\frac12.
\]
对 $t\ge0$ 定义\sourcefootnote{原书下式令 $f_0(t)\equiv a_0$；若
$a_0=f(0)\ne0$，则按此式构造的 $f$ 在 $t\ge1$ 并不消失，
与习题开头的目标不符。此处照录原式，不替作者选定补救步骤。}
\begin{equation}\label{eq:C-5-13}
  f(t):=\sum_{k=0}^\infty f_k(t),\qquad
  f_k(t):=a_k\left(\beta\left(\frac{t}{\varepsilon_k}\right)t\right)^k.
\end{equation}
证明 $f:\mathbb R\to\mathbb R$ 是光滑的。提示：对每个 $\ell\in\mathbb N$，
存在常数 $c_\ell>0$，使对所有 $t\ge0$ 与 $k\in\mathbb N$ 都有
\begin{equation}\label{eq:C-5-14}
  \bigl|f_k^{(\ell)}(t)\bigr|
  \le c_\ell|a_k|k^\ell\varepsilon_k^{k-\ell}.
\end{equation}
证明每个光滑\BookRecordIndex{T0088}殆复结构 $j\in\mathcal J(\mathbb D)$ 都可延拓为
$\mathbb C$ 上的光滑\BookRecordIndex{T0088}殆复结构，且在半径为二的圆盘外等于 $i$。
\end{exercise}
